\subsection{数据预处理：数据来源与字段结构}
原始计划来自附件 1 的 \texttt{Sheet1}，共有 151 行、5 列，其中首行为表头，实际计划为 150 条。每行对应一份用频计划，而不是一次单独使用事件。保留字段包括装备编号、频段区间、首次时间区间、间隔时长和使用次数。附件 2 的三个工作簿只承担问题一至问题三的结果输出，不被当作新的观测数据。

\begin{table}[H]
\centering
\caption{原始用频计划的数据结构与保留口径}
\label{tab:data-structure}
\begin{tabularx}{0.94\textwidth}{l l X}
\toprule
项目 & 数据规模或取值 & 处理口径 \\
\midrule
计划数量 & 150 条 & 编号唯一，按 A/B/C 三类分别为 20/40/90 条 \\
字段数量 & 5 列 & 全部保留，不覆盖原始附件 \\
频段宽度 & 3、10、15 & 保持原区间宽度，频段端点按整数处理 \\
首次时间宽度 & 2、3、5 & 作为单次使用时长 \(l_i\) \\
间隔时长 & 8、40、60 & 解释为一次占用结束到下一次占用开始之间的空闲间隔 \(g_i\) \\
使用次数 & 3、4、12 & 用于展开完整重复事件集合 \\
\bottomrule
\end{tabularx}
\end{table}

数据审计显示，五个字段没有空值、重复行或错误单元格；编号全部唯一，频段区间落在 \([0,100)\) 内。原始文件保持只读，所有展开结果和模型输入均写入新的结果目录。

\subsection{重复事件展开与端点处理}
设计划 \(i\) 的首次时间区间为 \([s_i,e_i)\)，单次时长为 \(l_i=e_i-s_i\)，空闲间隔为 \(g_i\)。第 \(k\) 次使用事件的时间区间定义为
\begin{equation}
I_{ik}=
\left[s_i+k(l_i+g_i),\ e_i+k(l_i+g_i)\right),
\qquad k=0,1,\ldots,n_i-1.
\label{eq:event-expansion}
\end{equation}
频段区间在重复事件之间保持不变。采用半开区间的好处是：一个事件在时刻 \(t\) 结束、另一个事件在同一时刻 \(t\) 开始时，两者不被判定为重叠。附件中的 A001 示例和由完整递推得到的 C083 末次事件分别用于核验展开语义和边界位置。

将每个事件映射到离散时频资源格 \(c=(t,f)\)，得到计划的占用集合 \(\mathcal K_i\)。连续半开区间相交判断用于主计算，离散资源格倒排索引用于独立复核。两种口径在问题一至问题三的验证结果中保持一致。

\subsection{跨问题模型输入}
预处理后的数据按问题形成三类模型输入：
\begin{enumerate}[label=(P\arabic*)]
  \item 问题一读取 150 份计划，展开得到 1,300 个重复事件，并构造装备对冲突清单；
  \item 问题二读取问题一的冲突背景，为每份计划枚举保留、频移、时移和撤销动作，形成 4,582 个候选动作及其资源格集合；
  \item 问题三读取经批准的问题二工作簿，固定 144 个活动计划后，在所有整数起点上枚举 52,136 个完整 C 类候选，再筛选得到 2,334 个与固定背景不冲突的候选。
\end{enumerate}

问题二和问题三的基准资源域为由完整数据展开得到的 \([0,643)\times[0,100)\)。其中 643 是数据驱动的最晚结束时刻，不是题面直接给出的全局上界；改变该边界时，相关模型和结果需要重新求解。
